\chapter{二分查找和答案空间}

二分查找的本质不是在数组里找数，而是在单调结构中找边界。只要能定义判断函数 \code{check(x)}，并且它在答案空间上单调，就可以二分。

基础二分在有序数组中查找目标值。写法看似简单，错误通常来自边界没有排除 \code{mid}、返回值不清楚或 \code{left + right} 溢出。

\begin{lstlisting}[language=C++,caption={基础二分查找}]
int binary_search_index(const std::vector<int>& nums, int target)
{
    int left = 0;
    int right = static_cast<int>(nums.size()) - 1;
    while (left <= right) {
        const int mid = left + (right - left) / 2;
        if (nums[mid] == target) {
            return mid;
        }
        if (nums[mid] < target) {
            left = mid + 1;
        } else {
            right = mid - 1;
        }
    }
    return -1;
}
\end{lstlisting}

找左边界时，标准库 \code{lower_bound} 是最好的参照：它找第一个不小于目标的位置。很多题不是找一个等于目标的位置，而是找第一个满足条件的位置。

答案二分把“位置”换成“答案值”。爱吃香蕉的珂珂中，答案是速度。速度越大越容易在规定时间内吃完，所以 \code{check(speed)} 从 false 变成 true。我们要找第一个 true。分割数组最大值、运输包裹能力、制作花束最少天数也都是这个模型。

\begin{principlebox}{答案二分三步}
第一，写出答案范围。第二，写出 \code{check(x)}。第三，证明单调性，并说明要找最小可行值还是最大可行值。没有单调性，不要二分。
\end{principlebox}

本章力扣主线：704 二分查找、35 搜索插入位置、34 查找第一个和最后一个位置、33 搜索旋转排序数组、162 寻找峰值、875 爱吃香蕉的珂珂、1011 在 D 天内送达包裹的能力、410 分割数组的最大值。
